% =====================================================================
%  2. The gradient: direction of steepest slope
% =====================================================================
\section{The gradient: direction of steepest slope}

\begin{frame}{From derivative to gradient}
  \begin{columns}[T]
    \begin{column}{0.5\textwidth}
      \small
      One variable: $J'(x)$ is the slope. Several: stack the partial derivative for \alert{each} variable,
      \[
        \nabla J(\x) = \begin{pmatrix}\partial J/\partial x_1\\ \vdots \\ \partial J/\partial x_n\end{pmatrix}.
      \]
      Example: $J = x_1^2 + x_2^2$ gives $\nabla J = (2x_1,\ 2x_2)$.

      \begin{keybox}[What the gradient tells us]
        \footnotesize At every point, $\nabla J$ is a vector that gives the \textbf{direction} and the \textbf{size} of the steepest
        slope. Moving along $-\nabla J$ goes downhill fastest.
      \end{keybox}
    \end{column}
    \begin{column}{0.48\textwidth}
      \centering
      \includegraphics[width=\linewidth,height=0.56\textheight,keepaspectratio]{fig_gradient_field}

      \footnotesize Arrows: $-\nabla J$ on $J = x_1^2 + x_2^2$, pointing towards the minimum.
    \end{column}
  \end{columns}
\end{frame}

\begin{frame}{Why $-\nabla J$ is the steepest way down}
  \begin{columns}[T]
    \begin{column}{0.55\textwidth}
      \small
      Take a small step $t$ in a unit direction $\vd$. By Taylor,
      $J(\x + t\vd) \approx J(\x) + t\,\nabla J(\x)\T\vd$, so the rate of change along $\vd$ is
      $\nabla J\T\vd$.
      \begin{thmbox}[Claim]
        Among all unit directions, $\nabla J\T\vd$ is smallest at $\vd^\star = -\nabla J/\norm{\nabla J}$,
        where it equals $-\norm{\nabla J}$.
      \end{thmbox}
      \textbf{Proof.} By Cauchy--Schwarz, $|\nabla J\T\vd|\le\norm{\nabla J}\norm{\vd} = \norm{\nabla J}$, so
      $\nabla J\T\vd\ge-\norm{\nabla J}$. Equality holds when $\vd$ parallel to $\nabla J$, pointing the opposite way. $\blacksquare$

      \vspace{3pt}
      \footnotesize This identifies the direction of steepest descent~\cite{boyd2004,cauchy1847}.
    \end{column}
    \begin{column}{0.43\textwidth}
      \centering
      % Steepest-descent geometry: the half-plane of descent directions, the unit
% circle of candidate directions d, the gradient and the angle theta.
\begin{tikzpicture}[scale=1.4, >={Stealth[length=7pt]}]
  % descent half-plane  {d : grad^T d < 0}  (gradient direction (0.447, 0.894))
  \begin{scope}
    \clip (-1.75,-1.35) rectangle (1.75,1.35);
    \fill[mGreen!9] (-2,1) -- (2,-1) -- (2,-3) -- (-2,-3) -- cycle;
  \end{scope}
  \draw[mGrey, dashed, line width=0.6pt] (-1.75,0.875) -- (1.75,-0.875);
  \node[font=\scriptsize, mGreen, anchor=south east] at (1.60,-1.23) {descent directions};
  \node[font=\scriptsize, mGrey, rotate=-26.6, anchor=south west] at (-1.72,0.94) {level set};
  % unit circle
  \draw[mDarkTeal!35, line width=1pt] (0,0) circle (1);
  % candidate directions (faint)
  \foreach \t in {200,230,260,290,320,350,80,110,140,170}
    \draw[mGrey!45, -{Stealth[length=4pt]}] (0,0) -- (\t:1);
  % a generic direction d and the angle
  \draw[->, line width=1.3pt, mBlue] (0,0) -- (-15:1) node[right, font=\small] {$\vd$};
  \draw[mBlue, line width=0.7pt] (63.4:0.43) arc[start angle=63.4, end angle=-15, radius=0.43];
  \node[mBlue, font=\small] at (24:0.67) {$\theta$};
  % gradient and steepest descent
  \draw[->, line width=1.8pt, mRed] (0,0) -- (63.4:1.18) node[above, font=\small] {$\nabla f(\x)$};
  \draw[->, line width=1.8pt, mGreen] (0,0) -- (243.4:1) node[below left=3pt, font=\small] {$\vd^\star$};
  \fill[mInk] (0,0) circle (1.4pt);
  \node[font=\small, mInk, fill=mPaper, inner sep=1pt] at (-0.28,0.27) {$\x$};
\end{tikzpicture}


      \footnotesize The rate along $\vd$ is $\norm{\nabla J}\cos\theta$, lowest at $\theta = 180^\circ$.
    \end{column}
  \end{columns}
\end{frame}

\begin{frame}{Contour plots, and why GD crosses them at right angles}
  \begin{columns}[T]
    \begin{column}{0.5\textwidth}
      \small
      A contour plot draws $J(x_1, x_2)$ as rings of equal height.

      \textbf{Why $\nabla J\perp$ contour.} Move along a contour by a small $\mathbf t$; the height does not change:
      \[
        J(\x + \mathbf t) - J(\x)\approx\nabla J\T\mathbf t = 0 .
      \]
      So each GD step leaves its contour at a right angle.
      \begin{keybox}[Reading a contour plot]
        \footnotesize Close contours: steeper slopes, larger steps. Wide gaps: smaller steps.
      \end{keybox}
    \end{column}
    \begin{column}{0.48\textwidth}
      \centering
      \includegraphics[width=\linewidth,height=0.7\textheight,keepaspectratio]{fig_gradient_field}
    \end{column}
  \end{columns}
\end{frame}

\begin{frame}{When the gradient is zero}
  \small
  Setting $\nabla J = \mathbf 0$ finds \emph{stationary} points. The second derivative tells them apart:

  \vspace{4pt}
  \centering
  \begin{tabular}{@{}l l l l@{}}
    \toprule
    Function & Stationary point & Second derivative & Type \\
    \midrule
    $J = x^2$ & $J' = 2x = 0\Rightarrow x = 0$ & $J'' = 2 > 0$ & \textcolor{mTeal}{\textbf{minimum}} \\
    $J = -x^2$ & $x = 0$ & $J'' = -2 < 0$ & \textcolor{mOrange}{\textbf{maximum}} \\
    $J = x^3$ & $J' = 3x^2 = 0\Rightarrow x = 0$ & $J'' = 6x = 0$, sign changes & \textcolor{mRed}{\textbf{saddle}} \\
    $J = x_1^2 + x_2^2$ & $\nabla J = (2x_1, 2x_2) = \mathbf 0$ & both curvatures $2 > 0$ & \textcolor{mTeal}{\textbf{minimum}} at $(0,0)$ \\
    \bottomrule
  \end{tabular}

  \vspace{6pt}
  \raggedright
  \begin{warnbox}[Zero gradient is not the same as a minimum]
    GD stops wherever the gradient vanishes, so it can stall at a saddle or settle in a \emph{local} minimum.
    Problem M1 shows this on the quartic $J = x^4 - 6x^3 + 8x^2$.
  \end{warnbox}
\end{frame}

\begin{frame}{What the three kinds of point look like}
  \begin{columns}[T]
    \begin{column}{0.5\textwidth}
      \centering
      \includegraphics[width=\linewidth, height=0.62\textheight, keepaspectratio]{fig_quartic_surface}
    \end{column}
    \begin{column}{0.48\textwidth}
      \centering
      \includegraphics[width=\linewidth, height=0.62\textheight, keepaspectratio]{fig_quartic_contour}
    \end{column}
  \end{columns}
  \vspace{2pt}
  \footnotesize $J = x_1^4 + x_2^4 - 4x_1x_2$ has $\nabla J = \mathbf 0$ at three points: two \textcolor{mTeal}{\textbf{minima}}
  at $\pm(1, 1)$ and a \textcolor{mRed}{\textbf{saddle}} at the origin. The arrows ($-\nabla J$) flow into the valleys,
  and at the saddle they arrive along one direction and leave along the other.
\end{frame}
